% ---------------------------------------------------------------------------
% Method. Formulas first, then the routing rule they imply, then the figure.
% ---------------------------------------------------------------------------
\section{Method}
\label{sec:method}

\subsection{Where the gradient is zero}
\label{sec:silence}

Let a \emph{unit} be a prompt with $G$ sampled answers and binary rewards
$r_1,\dots,r_G\in\{0,1\}$. GRPO centres rewards within the group,
%
\begin{equation}
  A_i \;=\; r_i - \bar r,
  \qquad
  \bar r \;=\; \frac{1}{G}\sum_{j=1}^{G} r_j ,
  \label{eq:adv}
\end{equation}
%
so if every sample in the group scores alike then $A_i=0$ for all $i$ and the unit
contributes exactly nothing to the update --- all $G$ sequences, every token. Writing
$p$ for the per-sample success probability, the probability that a unit carries any
signal at all is
%
\begin{equation}
  I_{\mathrm{RL}}(p,G) \;=\; 1 - p^{G} - (1-p)^{G},
  \label{eq:irl}
\end{equation}
%
which is zero exactly at $p=0$ and $p=1$ and is maximised at $p=\tfrac12$. Equation
\eqref{eq:irl} is not an approximation: it is the probability that the group is not
unanimous, and unanimity is precisely the condition under which \eqref{eq:adv} vanishes.

Two consequences drive the method.

\paragraph{The silent fraction is large and measurable.} Instrumenting a live run
(Sec.~\ref{sec:results}) puts $\hat{P}(\text{silent})$ between $0.29$ and $0.57$ depending on
the training step, at $64$ groups per batch: a third to a half of every batch is RL-dead.
Even the low end is an order of magnitude above the reach of the token-level shared-prefix
channel we measured first ($1.7\%$), and that gap explains, without appealing to noise, why
an intervention confined to that channel moved nothing.

\paragraph{The two silent regimes need opposite responses, and they are not equally
common.} A unit silent because $\hat p = 1$ and a unit silent because $\hat p = 0$ are both
invisible to \eqref{eq:adv}, but they are not the same object, and routing them identically
wastes one of them. Measured on the control arm, the silent channel is $87.5\%$ solved
(Table~\ref{tab:split}), which decides which of the two branches carries the method.

\paragraph{The composition is a property of the task \emph{and} the model, and it moves in both
directions.} The $87.5\%$ is GSM8K with Qwen2.5-1.5B-Instruct; on MATH, holding the model
fixed, it inverts to $39.1\%$ solved and $60.9\%$ unsolved (Table~\ref{tab:flip}). Re-measured
on the arm this paper trains --- decontaminated DeepMath with Qwen2.5-32B-Instruct, over $245$
logged steps of the baseline --- the channel is $0.505$ of all groups and splits $0.338$ solved,
$0.167$ unsolved and $0.000$ unclassified, i.e.\ $66.9\%$ solved and $33.1\%$ unsolved. The
unsolved share is therefore roughly \emph{half} what the MATH column reports, and we state the
confound plainly: that pair varies the model as well as the task, so it cannot separate a
task effect from a scale effect, and both readings shrink the branch. Two consequences are load
bearing. First, any statement sized on the $60.9\%$ figure must not be carried onto this arm:
the unsolved branch is about a sixth of every batch here, not a third. Second, the
$0.000$ unclassified mass means no group on this arm is silent for a reason the correctness
split cannot express, so unlike on GSM8K the composition ratio has a clean denominator and
$\text{unsolved}/\text{silent}$ and $\text{unsolved}/(\text{solved}+\text{unsolved})$ coincide.

\subsection{A free target exists on exactly the solved half}
\label{sec:selftarget}

The asymmetry is sharper than "different regimes''. For a unit with $\hat p > 0$, at least
one sampled answer was correct, and that answer \emph{is} a supervised target drawn from
the model's own rollout. No external teacher is required; this is rejection-sampling
fine-tuning. Define
%
\begin{equation}
  \mathrm{selfTarget}(u) \;=\; \mathbb{1}\!\left[\hat p_u > 0\right],
  \qquad
  \mathrm{target}(u) \;=\; \mathrm{selfTarget}(u) \;\vee\; \mathrm{teacher}(u).
  \label{eq:target}
\end{equation}
%
The units on which $I_{\mathrm{RL}}=0$ because $\hat p = 1$ are therefore exactly the units
on which a target is \emph{guaranteed} to exist at zero marginal cost. The units on which
$I_{\mathrm{RL}}=0$ because $\hat p = 0$ are exactly the units on which no target exists
unless one is supplied. The gradient-free half of the batch splits cleanly into a half that
is free to reuse and a half that is not reusable at all.

\subsection{Supplying a target for the half that has none}
\label{sec:supply}

``Not reusable at all'' is a statement about \eqref{eq:adv}, not about the batch. An advantage
is a coefficient on tokens the model \emph{emitted}, so on a unit with $\hat p = 0$ every value
that could be written into the advantage tensor multiplies a wrong derivation, and a positive
constant there is a step toward the error --- which is why $c_- \le 0$ is required and why no
router in our registry reaches this branch. The only altitude at which a correct row can receive
gradient is \emph{batch construction}: put the correct response in as a row, with its own prompt
and its own loss mask, and the estimator that already runs on every other row runs on it too.
This is the construction of LUFFY and of DyME, and it is what makes the unsolved branch
addressable at all.

\paragraph{The supplier is an axis, not a constant.} Once the row enters by construction, where
it came from is a separate question from where it goes: the splice, the masks, the counters and
the off-policy treatment are properties of inserting a correct row and do not depend on its
origin. We therefore define a supplier interface with four implementations --- the dataset's
gold solution; a \emph{self-generated} correct rollout of the same prompt from an earlier step
or a higher-temperature resample; a row from an offline \emph{corpus} of solved examples; and a
verified completion from a stronger \emph{teacher} --- and route a unit to one of them. Prompt
identity is derived from the prompt tokens rather than from a batch-local index, because a
self-generated row is by definition one the same prompt received at a different step.

\paragraph{A supplier that has nothing must refuse, and the refusal is the measurement.} Each
supplier either returns token ids or raises a typed refusal, and every reason has its own
counter that is reported whether or not it is zero: a missing dataset solution, a prompt never
solved before, a derivation too long for the row, and a teacher completion that fails
verification have four different remedies, and a single ``skipped'' total sends the reader to
the wrong one. No supplier substitutes a row belonging to a different prompt, and no supplier
truncates: a cut-off derivation is a wrong target that still looks like a target. The teacher
requires an explicit verifier for the same reason --- a stronger model is not a correct one, and
this is precisely the branch on which nothing else would catch the error.

\paragraph{The source decision is fixed, and its control is not optional.} The choice of
supplier is expressible per unit, but we do not learn it here. We use an explicit ordered
policy, cheapest-and-most-trusted first, and compare it against a control that replays the
policy's own realised source assignment permuted across units, so the mixture of sources is
matched exactly by construction and only the \emph{targeting} differs. This is the same
falsification the routing control performs one axis over, and we run it for the same reason: in
this project a learned mechanism has repeatedly matched a feature-blind control at matched
proportions once the control was finally run, and a claim about choosing sources that has not
been separated from a claim about mixing them is not a claim.

\subsection{Routing: the harness is a destination, not a schedule}
\label{sec:routing}

Prior work treats the choice of what to evolve as a property of the \emph{run}: a schedule
or a global flag. We make it a property of the \emph{unit}, and we make it two-dimensional.
A routing decision is a pair
%
\begin{equation}
  \pi(u) \;=\; \bigl(m(u),\, h(u)\bigr),
  \qquad
  m \in \{\textsc{rl}, \textsc{sft}, \textsc{skip}\},
  \qquad
  h \in \{\textsc{none}, \textsc{propose}, \textsc{validate}\},
  \label{eq:action}
\end{equation}
%
where $m$ names the estimator applied to the unit's tokens and $h$ names what the unit
contributes to harness evolution. The two axes are orthogonal by construction, so a single
trajectory can feed the gradient and the harness at once. This is deliberate: the two
consumers read different things from a trajectory --- the estimator reads the tokens, the
harness evolver reads the failure mode --- and a partition discards one of the readings.
Co-Harness's success$\to$model / failure$\to$harness rule is the special case in which
$\pi$ is constrained to be a partition, and we retain it as an ablation
(\texttt{partition=True}) rather than a rewrite.

The rule implied by \eqref{eq:irl} and \eqref{eq:target} is then forced rather than chosen:
%
\begin{equation}
  \pi(u) \;=\;
  \begin{cases}
    (\textsc{sft}_{\text{self}},\ \textsc{validate})
      & \hat p_u = 1 \quad\text{(RL dead; target free)}\\[2pt]
    (\textsc{sft}_{\text{teacher}},\ \textsc{propose})
      & \hat p_u = 0 \ \wedge\ \mathrm{teacher}(u) \\[2pt]
    (\textsc{skip},\ \textsc{propose})
      & \hat p_u = 0 \ \wedge\ \neg\,\mathrm{teacher}(u)
      \quad\text{(harness is the only consumer)}\\[2pt]
    (\textsc{rl},\ \textsc{none})
      & 0 < \hat p_u < 1 \quad\text{(RL has signal; leave it alone)}
  \end{cases}
  \label{eq:rule}
\end{equation}
%
Note the third case. When a unit is unsolved and no teacher is available, the model arm has
nothing to learn from it, and any method that only evolves weights must discard the unit.
Routing it to the harness is what makes it recoverable at all, which is the concrete sense
in which the harness axis is not a convenience.

\paragraph{One trajectory, two consumers.} \eqref{eq:rule} as written still couples the two
axes: a unit that RL can learn from contributes nothing to the harness. That is a residual
partition, and it discards information --- a \emph{mixed} unit contains failed samples, and
their failure mode is exactly what a harness evolver reads, while the RL update consumes only
the reward. We therefore decouple the harness action with its own threshold $\tau_h$:
%
\begin{equation}
  h(u) \;=\;
  \begin{cases}
    \textsc{validate} & \hat p_u = 1 \\
    \textsc{propose}  & \hat p_u \le \tau_h \\
    \textsc{none}     & \text{otherwise,}
  \end{cases}
  \qquad \tau_h \ge \tau_{\text{fail}} ,
  \label{eq:harness}
\end{equation}
%
so at $\tau_h = \tau_{\text{fail}}$ the rule is unchanged, and raising it lets a unit take
$(\textsc{rl}, \textsc{propose})$ --- feeding the gradient and the harness from the same
trajectory. Under $\texttt{partition}$ this is disabled, so the Co-Harness special case is
preserved exactly.

\paragraph{Rollback.} Every capability above is gated. With the harness arm disabled,
$h(u) \equiv \textsc{none}$; with routing disabled, $\pi(u) \equiv (\textsc{rl},
\textsc{none})$ for every unit and the update is bit-identical to unmodified GRPO; with the
supplier axis of Sec.~\ref{sec:supply} disabled, the batch construction reduces to its single
original source and its output is byte-identical, key set included, which we verify against a
digest recorded before the generalisation existed rather than by inspection. Each
ablation in Sec.~\ref{sec:results} is one flag away from the vanilla baseline, so a
difference cannot be attributed to an incidental change of configuration.



\subsection{Who decides: from a threshold to a controller}
\label{sec:controller}

Everything above specifies an action space. What chooses within it is a separate question,
and it is where the contribution lies --- a rule keyed on $\hat p$ is a heuristic, and a
learned policy that sees only $\hat p$ cannot beat the best threshold on $\hat p$ by more
than noise. We therefore give the controller features, and instantiate it three ways so the
comparison is an ablation rather than a demonstration.

\paragraph{Features, at zero marginal cost.} Each unit is described by seven statistics
computed from tensors the batch \emph{already carries} --- rewards, the loss mask, and the
sampled log-probabilities. No extra forward pass, no extra generation. The constraint is
deliberate: a feature set needing its own inference budget would have to beat spending that
budget on more rollouts, which is a far higher bar than beating a threshold. Two of the seven
carry information $\hat p$ cannot: the \emph{truncated fraction} separates a unit that failed
from one that merely ran out of budget, and the \emph{log-probability dispersion} separates a
group unanimous in its answer from one also unanimous in its reasoning. A solve-rate
threshold collapses both distinctions by construction.

\paragraph{Three controllers, one action space.}
%
\begin{itemize}
  \item \textbf{Rule} --- thresholds on $\hat p$, as in \eqref{eq:rule}. The baseline the
        others must beat, not a strawman: it encodes the exact zero-gradient condition.
  \item \textbf{Learned weights} --- LinUCB, one ridge model per mode over the features,
        updated from observed outcomes. Ridge rather than a network because the outcome
        channel is confounded and low-signal --- one scalar per batch, produced by every
        decision in it --- and capacity would mostly fit that noise.
  \item \textbf{Learned code} --- the policy is generated \emph{source}, validated against an
        allowlist and cost-vetted in a subprocess before adoption. Strictly larger
        hypothesis class than weights: it can compose features into ratios and nested cases
        rather than only reweighting them.
\end{itemize}
%
\paragraph{The control that decides the claim.} A fourth arm shuffles which units receive
which mode while holding the \emph{proportion} of units in each mode fixed. Without it, a
controller that helps cannot be distinguished from a change in how much SFT the run did ---
and the latter is not a controller at all. Any improvement reported for the learned arms is
the improvement over this control, not over the fixed rule alone.

\paragraph{What we verify about the learned arm, and what we do not.} On a synthetic
environment whose two regions differ \emph{only} in the truncated fraction, with $\hat p$
held identical, the contextual controller reaches the optimal mode in both regions across
seeds, while a context-free bandit on the same stream provably cannot exceed $0.5$ in the
worse region. That establishes the mechanism can use a feature a threshold cannot see. It
does \emph{not} establish that it learns under the confounded batch-level channel of a real
run, which is a separate claim and is not made here.

\subsection{The procedure, end to end}
\label{sec:algo}

Algorithm~\ref{alg:route} is the whole method. Every line is either an existing GRPO step or
a decision from \eqref{eq:rule}--\eqref{eq:harness}; nothing else is added, and with routing
disabled lines 6--13 do not execute.

\begin{algorithm}[t]
\caption{Per-batch signal routing. $\mathcal{B}$ is a batch of prompts, $G$ samples each.}
\label{alg:route}
\begin{algorithmic}[1]
\Require thresholds $\tau_{\text{solve}}{=}1$, $\tau_{\text{fail}}{=}0$, $\tau_h \ge \tau_{\text{fail}}$;
         weights $c_+ \ge 0$, $c_- \le 0$; flags \texttt{routing}, \texttt{harness}
\State roll out $G$ samples per prompt; score to binary rewards $r_{u,i}$
\State $\hat p_u \gets \frac{1}{G}\sum_i r_{u,i}$ \Comment{group solve rate}
\State $A_{u,i} \gets r_{u,i} - \hat p_u$ \Comment{group-centred advantages, \eqref{eq:adv}}
\State $\mathcal{S} \gets \{u : \max_i |A_{u,i}| = 0\}$ \Comment{RL-silent: exactly the unanimous units}
\If{\texttt{routing}}
  \ForAll{$u \in \mathcal{S}$}
    \If{$\hat p_u \ge \tau_{\text{solve}}$} \Comment{solved: a correct sample is its own target}
      \State $A_{u,i} \gets A_{u,i} + c_+$ for all response tokens
    \ElsIf{$\hat p_u \le \tau_{\text{fail}}$} \Comment{failed: no self-target exists}
      \State $A_{u,i} \gets A_{u,i} + c_-$ \textbf{if} $c_- \neq 0$, else leave at $0$
    \EndIf
  \EndFor
\EndIf
\If{\texttt{harness}}
  \State $\mathcal{P} \gets \{u : \hat p_u \le \tau_h\}$; \; $\mathcal{V} \gets \{u : \hat p_u = 1\}$
  \State emit $(\textsc{propose}, u)$ for $u \in \mathcal{P}$ and $(\textsc{validate}, u)$ for $u \in \mathcal{V}$
\EndIf
\State PPO update on $\{A_{u,i}\}$, clipping unchanged
\end{algorithmic}
\end{algorithm}

\paragraph{Why the update needs no new loss term.} Line 8 is the whole of the supervised
branch. For a solved unit every $A_{u,i}$ is zero, so adding $c_+$ makes the objective
$c_+\sum_i \log \pi(y_{u,i})$ over samples already known to be correct --- which is supervised
fine-tuning on them, with the importance ratio at $1$ because the data is on-policy. PPO's
clip still bounds the step, which matters because sharpening an already-correct policy is the
mechanism's predicted failure. Line 10 is its mirror: $c_- \le 0$ pushes down samples known to
be wrong, and $c_- > 0$ is rejected rather than trusted, since it would train the model to
reproduce them.

\paragraph{Cost.} Lines 5--13 add no generation. The batch is the batch that GRPO already
rolled out; the silent units are re-read, not re-sampled. This is the whole of the difference
from DAPO, which reaches the same set by discarding it and paying to generate a replacement.

% Method flowchart. Standalone-compatible: the same TikZ source is exported to SVG by
% scripts/build_figs.sh, so the figure in the paper and the figure on the web are one file.
\begin{figure}[t]
\centering
\begin{tikzpicture}[
  font=\small,
  node distance=6mm and 9mm,
  box/.style   ={rectangle, rounded corners=2pt, draw=black!55, fill=black!3,
                 minimum height=7.5mm, inner xsep=3mm, align=center},
  test/.style  ={diamond, aspect=2.1, draw=black!55, fill=black!4,
                 inner xsep=0.4mm, inner ysep=0mm, align=center},
  sink/.style  ={rectangle, rounded corners=2pt, draw=black!70, fill=black!8,
                 minimum height=7.5mm, inner xsep=3mm, align=center},
  arr/.style   ={-{Latex[length=2mm]}, draw=black!65},
  lbl/.style   ={font=\scriptsize, inner sep=1pt, fill=white}
]

\node[box] (roll) {rollout: $G$ samples for unit $u$};
\node[box, below=of roll] (rate) {$\hat p_u=\frac1G\sum_i r_i$};
\node[test, below=of rate] (sil) {$I_{\mathrm{RL}}(\hat p_u,G)=0$?\\[-1pt]\scriptsize(unanimous)};

\node[sink, right=22mm of sil] (rl) {$(\textsc{rl},\ \textsc{none})$\\[-1pt]\scriptsize GRPO as usual};

\node[test, below=9mm of sil] (solved) {$\hat p_u = 1$?};
\node[sink, right=22mm of solved] (sft) {$(\textsc{sft}_{\text{self}},\ \textsc{validate})$\\[-1pt]\scriptsize target = own correct sample};

\node[test, below=9mm of solved] (tea) {teacher\\available?};
\node[sink, right=22mm of tea] (sftt) {$(\textsc{sft}_{\text{teacher}},\ \textsc{propose})$};
\node[sink, below=9mm of tea] (skip) {$(\textsc{skip},\ \textsc{propose})$\\[-1pt]\scriptsize harness is the only consumer};

\draw[arr] (roll) -- (rate);
\draw[arr] (rate) -- (sil);
\draw[arr] (sil)   -- node[lbl,above]{no} (rl);
\draw[arr] (sil)   -- node[lbl,left]{yes} (solved);
\draw[arr] (solved)-- node[lbl,above]{yes} (sft);
\draw[arr] (solved)-- node[lbl,left]{no ($\hat p_u{=}0$)} (tea);
\draw[arr] (tea)   -- node[lbl,above]{yes} (sftt);
\draw[arr] (tea)   -- node[lbl,left]{no} (skip);

\begin{scope}[on background layer]
  \node[draw=black!25, dashed, rounded corners=3pt, fit=(solved)(tea)(skip)(sftt)(sft),
        inner sep=3mm, label={[font=\scriptsize\itshape, black!55]above right:%
        RL-silent: 29--57\% of groups, measured}] {};
\end{scope}
\end{tikzpicture}
\caption{%
  \textbf{Routing decides, per unit, both an estimator and a harness contribution.}
  The branch on $I_{\mathrm{RL}}$ is exact, not heuristic: a unanimous group has every
  advantage identically zero by \eqref{eq:adv}. The dashed region is the part of the batch
  ordinary GRPO cannot learn from, measured at $29$--$57\%$ of groups depending on training
  step. It does not split evenly: $87.5\%$ of it is solved, where a supervised target is
  free (the unit's own correct sample), and the remainder has no target unless one is
  supplied --- and on that remainder the harness is the only consumer that can use the unit
  at all.}
\label{fig:routing}
\end{figure}

